\par\Needspace{6\baselineskip}\medskip
\begin{center}\hypertarget{native-225-article162}{}\pdfbookmark[2]{162.}{article-162}162.\end{center}
In casu \emph{secundo}, ubi $B''$ cum $B$ coincidit, etiam $D'$ cum iisdem vel
cum puncto opposito coincidet. Erunt itaque $AD'-\delta$ et $A''D'-\delta''$
vel $=0$ vel $=180^0$: unde ex aequationibus art. 143 derivamus
\SourceLeaf{226}\[
\frac{n'r'}{n}=\pm\frac{\sin\epsilon'}{\sin\epsilon}\cdot
\frac{R\sin\delta}{\sin(z+A'D-\delta')}
\]
\[
\frac{n'r'}{n''}=\pm\frac{\sin\epsilon'}{\sin\epsilon''}\cdot
\frac{R''\sin\delta''}{\sin(z+A'D''-\delta')}
\]
\[
\begin{aligned}
&R\sin\delta\sin\epsilon''\sin(z+A'D''-\delta')\notag\\
&\qquad=PR''\sin\delta''\sin\epsilon\sin(z+A'D-\delta').
\end{aligned}
\]
Hinc manifestum est, $z$, independenter a $Q$, per solam $P$ determinabilem esse
(nisi forte fuerit $A'D''=A'D$ vel $=A'D\pm180^0$, ubi ad casum tertium
delaberemur): inventa autem $z$, innotescet etiam $r'$, et proin adiumento
valorum quantitatum $\dfrac{n'r'}{n}\mathpunct{\raisebox{.45ex}{,}}$
$\dfrac{n'r'}{n''}$ etiam $\dfrac{n}{n'}$ et
$\dfrac{n''}{n'}\mathpunct{\raisebox{.45ex}{;}}$ hinc denique etiam
$Q=2\left(\dfrac{n}{n'}+\dfrac{n''}{n'}-1\right)r'^3$.
Manifesto igitur $P$ et $Q$ tamquam data ab invicem independentia considerari
nequeunt, sed vel unicum tantummodo datum exhibebunt, vel data incongrua.
Situs punctorum $C$, $C''$ in hoc casu arbitrarius manebit, si modo in eodem
circulo maximo cum $C'$ capiantur.

In casu \emph{tertio}, ubi $A'$, $B$, $B'$, $B''$ in eodem circulo maximo
iacent, $D$ et $D''$ resp. cum punctis $B''$, $B$, vel cum punctis oppositis
coincident: hinc e combinatione aequationum VII, VIII, IX art. 143 colligitur
\[
P=\frac{R\sin\delta\sin\epsilon''}{R''\sin\delta''\sin\epsilon}
=\frac{R\sin(l'-l)}{R''\sin(l''-l')}\mathpunct{\cdot}
\]
In hoc itaque casu valor ipsius $P$, per ipsa problematis data iam habetur,
adeoque positio punctorum $C$, $C'$, $C''$ indeterminata manebit.
